\documentclass[12pt,a4paper]{article}
\usepackage[utf8]{inputenc}
\usepackage[T1]{fontenc}
\usepackage[brazil]{babel}
\usepackage{amsmath,amssymb}
\usepackage{geometry}
\geometry{margin=2.5cm}
\title{Material de Estudo: Prepara\c{c}\~ao para Prova I}
\author{Princ\'ipio de Modelagem Matem\'atica}
\date{Unidade 15}
\begin{document}
\maketitle

\section*{Objetivo}
Consolidar todos os t\'opicos das Unidades 1--14 com um guia de f\'ormulas essenciais, uma lista de verificação e exerc\'icios-modelo.

\section*{Guia de F\'ormulas Essenciais}

\subsection*{An\'alise Dimensional}
\[
\text{Buckingham: } n \text{ vari\'aveis}, k \text{ dimens\~oes} \Rightarrow n-k \text{ grupos adimensionais.}
\]
\[
Re = \rho v L/\mu, \quad Fr = v/\sqrt{gL}, \quad Ma = v/c_\text{som}.
\]

\subsection*{Escalas e Leis de Pot\^encia}
\[
Y = aX^b. \quad \text{Kleiber: } B \propto M^{3/4}.
\]

\subsection*{Fractais}
\[
d_f = \frac{\log N}{\log r}. \quad \text{Koch: } d_f = \log4/\log3 \approx 1{,}26.
\]

\subsection*{Taylor}
\[
f(x) = \sum_{n=0}^\infty \frac{f^{(n)}(a)}{n!}(x-a)^n. \quad |R_n| \le \frac{M|x-a|^{n+1}}{(n+1)!}.
\]

\subsection*{Estat\'istica e Propaga\c{c}\~ao}
\[
\bar x = \frac{1}{N}\sum x_i, \quad s = \sqrt{\frac{\sum(x_i-\bar x)^2}{N-1}}, \quad s_{\bar x} = \frac{s}{\sqrt{N}}.
\]
\[
\sigma_z^2 = \left(\frac{\partial f}{\partial x}\right)^2\!\sigma_x^2 + \left(\frac{\partial f}{\partial y}\right)^2\!\sigma_y^2.
\]

\subsection*{Regress\~ao}
\[
\hat\beta_1 = \frac{\sum(x_i-\bar x)(y_i-\bar y)}{\sum(x_i-\bar x)^2}, \quad R^2 = 1-\frac{SS_\text{res}}{SS_\text{tot}}.
\]

\subsection*{EDO e Crescimento}
\[
\frac{dx}{dt}=kx \Rightarrow x(t)=x_0 e^{kt}. \quad t_{1/2}=\ln2/|k|.
\]
\[
\text{Log\'istico: }N(t) = \frac{K}{1+(K/N_0-1)e^{-rt}}.
\]

\subsection*{Tr\'afego}
\[
q=\rho v. \quad \text{Greenshields: }q=v_f\rho(1-\rho/\rho_\text{max}). \quad q_\text{max}=v_f\rho_\text{max}/4.
\]
\[
\text{Choque: }w = \Delta q/\Delta\rho.
\]

\section*{Checklist de Revis\~ao}
\begin{itemize}
  \item[$\square$] Sei aplicar o teorema de Buckingham a um novo problema.
  \item[$\square$] Sei calcular a dimens\~ao fractal de um fractal auto-similar.
  \item[$\square$] Sei expandir $e^x$, $\sin x$ e $\cos x$ em s\'erie de Taylor e estimar o erro.
  \item[$\square$] Sei propagar incertezas atrav\'es de uma f\'ormula.
  \item[$\square$] Sei ajustar uma reta por m\'inimos quadrados e calcular $R^2$.
  \item[$\square$] Sei resolver EDOs de 1\textordfeminine\ ordem por separa\c{c}\~ao de vari\'aveis.
  \item[$\square$] Sei calcular velocidade de onda de choque em tr\'afego (Rankine-Hugoniot).
\end{itemize}

\section*{Exerc\'icios-Modelo de Prova}

\begin{enumerate}
  \item \textbf{(4 pts)} A for\c{c}a de arrasto $F_D$ depende de $\rho_f$ (densidade do flu\'ido), $v$ (velocidade) e $D$ (di\^ametro). Use o teorema $\Pi$ para deduzir $F_D = C_D \rho_f v^2 D^2$. Identifique o papel do coeficiente $C_D$.
  
  \item \textbf{(3 pts)} Calcule $\ln(1{,}05)$ usando a s\'erie de Taylor at\'e 3 termos. Qual o erro percentual em rela\c{c}\~ao ao valor exato?
  
  \item \textbf{(3 pts)} Dados $\{(1,2), (2,4{,}1), (3,5{,}9), (4,8{,}2)\}$, ajuste uma reta por OLS e calcule $R^2$.
  
  \item \textbf{(3 pts)} Uma subst\^ancia radioativa tem meia-vida de 30 anos. Que porcentagem resta ap\'os 90 anos? Ap\'os 100 anos?
  
  \item \textbf{(4 pts)} Numa rodovia com $v_f=90$ km/h e $\rho_\text{max}=180$ ve\'ic/km: (a) calcule $q_\text{max}$; (b) se $\rho_L=45$ ve\'ic/km e $\rho_R=135$ ve\'ic/km, calcule a velocidade da onda de choque; (c) interprete o sinal de $w$.
  
  \item \textbf{(3 pts)} Um objeto fractal tem $N=7$ c\'opias com fator de escala $r=4$. Calcule $d_f$ e interprete: isso \'e mais parecido com uma linha ou com um plano?
\end{enumerate}

\section*{Refer\^encias}
\begin{itemize}
  \item Material das Unidades 1--14.
  \item Notas de aula da disciplina.
\end{itemize}

\end{document}
